\documentclass[11pt,a4paper]{article}
\usepackage[utf8]{inputenc}
\usepackage[T1]{fontenc}
\usepackage[spanish,es-noshorthands]{babel}
\usepackage{textcomp}
\usepackage[margin=2.2cm]{geometry}
\usepackage{enumitem}
\usepackage{longtable}
\usepackage{titlesec}
\usepackage[hidelinks]{hyperref}
\setlist[itemize]{topsep=2pt,itemsep=1pt,parsep=0pt}
\setlist[enumerate]{topsep=2pt,itemsep=1pt,parsep=0pt}
\titleformat{\section}{\Large\bfseries}{}{0pt}{}
\titleformat{\subsection}{\large\bfseries}{}{0pt}{}
\titleformat{\subsubsection}{\normalsize\bfseries}{}{0pt}{}
\setlength{\parindent}{0pt}
\setlength{\parskip}{4pt}
\title{Laboratorio de guitarra\\[2pt]\large Material archivado: improvisaci\'on, melod\'ia y fuentes}
\author{}
\date{}
\begin{document}
\maketitle
Este documento guarda el contenido de los m\'odulos de Improvisaci\'on, Melod\'ia y Fuentes que se retiraron de la web para dejarla solo con Escalas. La fuente original era \texttt{js/content.js}.
\tableofcontents
\newpage
\section{T\'ecnicas de improvisaci\'on}
14 ejercicios ordenados de menos a m\'as, con su fuente citada.
\subsection{1. Las 4 formas de conocer el patrón}
\textbf{Base} \hfill \emph{fuente: Absolutely Understand Guitar \textemdash{} Scotty West (2001)}
\par Una escala se sabe de verdad cuando la podés pensar, leer, escuchar y cantar. Tocar dedos sin oreja es "tipear sin saber qué dice el texto".
\begin{enumerate}
\item Pensar: nombrá en voz alta la tónica y los grados (1 2 3 4 5 6 7) antes de tocar.
\item Leer: tocá mirando el diagrama, sin mirar la mano.
\item Escuchar: tocá y cantá la nota que acabás de tocar.
\item Cantar: cantá la escala completa y después tocala. Si no la podés cantar, no la tenés.
\end{enumerate}
\emph{Preset: escala major, 60 BPM.}
\subsection{2. Dos octavas subiendo y bajando}
\textbf{Técnica} \hfill \emph{fuente: Absolutely Understand Guitar \textemdash{} Scotty West (2001)}
\par La escala tiene que salir igual de cómoda en la subida y en la bajada, con metrónomo y sin apurar. Es el "trabajo de hormiga" que después libera.
\begin{enumerate}
\item Elegí una posición (box) y tocá dos octavas subiendo y bajando.
\item 2 notas por click, después 3 y después 4. Cambiá a 4 notas por click solo si suena parejo.
\item Pasá por las 5-7 posiciones sin cambiar de tempo.
\item Grabate 30 segundos y escuchá: ¿se nota dónde está el pulgar o dónde dudás?
\end{enumerate}
\emph{Preset: escala pent-minor, 60 BPM.}
\subsection{3. Secuencias numéricas dentro de la escala}
\textbf{Creatividad} \hfill \emph{fuente: Absolutely Understand Guitar \textemdash{} Scotty West (2001)}
\par Partir de números (no de sonidos) te obliga a escuchar. Además genera material melódico real: muchas líneas conocidas son secuencias.
\begin{enumerate}
\item Tocá 1-2-3-4, 2-3-4-5, 3-4-5-6... en 16avos sobre un solo box.
\item Probá 1-3-5-3, 2-4-6-4, 3-5-7-5.
\item Probá 1-2-3-1, 2-3-4-2 (célula que repite).
\item Ahora tocá la misma secuencia pero cambiando la figuración: 8vos, corcheas con swing, tresillos.
\end{enumerate}
\emph{Preset: escala major, 70 BPM.}
\subsection{4. Escalas por intervalos (3\textordfeminine{}, 4\textordfeminine{}, 5\textordfeminine{}, 6\textordfeminine{})}
\textbf{Oído} \hfill \emph{fuente: Absolutely Understand Guitar \textemdash{} Scotty West (2001)}
\par Las escalas están hechas de 2das, pero la música se mueve por saltos. Practicar por intervalos entrena el intervalo, no la posición.
\begin{enumerate}
\item Tocá la escala en 3\textordfeminine{}s: 1-3, 2-4, 3-5... (suena "vacío", es normal).
\item Pasá a 4\textordfeminine{}s, 5\textordfeminine{}s y 6\textordfeminine{}s. Cada intervalo tiene su carácter.
\item Cantá el intervalo antes de tocarlo.
\item Terminá eligiendo un intervalo y construí una frase entera solo con ese salto.
\end{enumerate}
\emph{Preset: escala major, 60 BPM.}
\subsection{5. Arpegios del campo armónico}
\textbf{Armonía} \hfill \emph{fuente: Absolutely Understand Guitar \textemdash{} Scotty West (2001)}
\par Cada grado de la escala mayor genera un acorde. Tocar sus tríadas es el puente entre escala y canción: los acordes salen de la misma escala.
\begin{enumerate}
\item Tocá las 7 tríadas diatónicas en un box: I, ii, iii, IV, V, vi, vii\textdegree{}.
\item Nombralas en voz alta con su calidad (mayor, menor, disminuido).
\item Tocá la tríada y después la escala del mismo grado. Vas a escuchar qué notas "pertenecen".
\item Inventá una frase terminando en la 3\textordfeminine{} del acorde que suena (nota objetivo).
\end{enumerate}
\emph{Preset: escala major, 64 BPM, fondo diatonic-major.}
\subsection{6. Notas de reposo vs. notas de paso}
\textbf{Improvisación} \hfill \emph{fuente: Absolutely Understand Guitar \textemdash{} Scotty West (2001)}
\par Con el mismo set de notas, la diferencia entre sonar bien y sonar mal es la PITCH y la DURACIÓN. El 4\textordmasculine{} y el 7\textordmasculine{} agregan color si pasan; si te quedás dormido en ellos, suenan a error.
\begin{enumerate}
\item Improvisá 1 minuto aterrizando SOLO en 1-3-5 (notas del acorde).
\item Ahora agregá el 4 y el 7, pero en corcheas cortas que van a resolver al 3 o al 1.
\item Probá quedarte 2 tiempos en el 4 sobre un acorde mayor: sentí la tensión. Resolve.
\item Regla: notas de reposo = duración larga; notas de paso = duración corta.
\end{enumerate}
\emph{Preset: escala dorian, 72 BPM, fondo drone.}
\subsection{7. Modos: el mismo set, distinto "Do"}
\textbf{Modos} \hfill \emph{fuente: Absolutely Understand Guitar \textemdash{} Scotty West (2001)}
\par Los 7 modos son el mismo grupo de notas: lo que cambia es dónde está la tónica. El modo no lo define la escala, lo define el acorde al que le tocás encima.
\begin{enumerate}
\item Tocá C jónico sobre un drone de C y después A eólico sobre un drone de A.
\item Fijate que son las 7 teclas blancas, pero suenan completamente distintas.
\item Escala por modo + drone de la tónica, 2 minutos cada uno.
\item Cantá la nota característica de cada modo: jónico (7), dórico (6 natural), frigio (\ensuremath{\flat}2), lidio (\ensuremath{\sharp}4), mixolidio (\ensuremath{\flat}7), eólico (\ensuremath{\flat}6), locrio (\ensuremath{\flat}5).
\end{enumerate}
\emph{Preset: escala lydian, 66 BPM, fondo drone.}
\subsection{8. Pentatónica + \ensuremath{\flat}5 = blues}
\textbf{Blues} \hfill \emph{fuente: Absolutely Understand Guitar \textemdash{} Scotty West (2001)}
\par La pentatónica menor es la caja más segura del rock y el blues, pero el \ensuremath{\flat}5 (blue note) agrega la tensión que la hace sonar a blues y no a escala.
\begin{enumerate}
\item Caja 1 de la pentatónica menor, dos octavas, con backing de blues.
\item Agregá el \ensuremath{\flat}5 solo como nota de paso entre el 4 y el 5.
\item Tocá la misma frase tres veces: sin \ensuremath{\flat}5, con \ensuremath{\flat}5 de paso, con \ensuremath{\flat}5 sostenido (sentí cómo cambia de sentido).
\item Probá llamada y respuesta: frase de 2 compases, silencio de 2 compases.
\end{enumerate}
\emph{Preset: escala blues, 84 BPM, fondo blues.}
\subsection{9. Pregunta y respuesta (call \& response)}
\textbf{Fraseo} \hfill \emph{fuente: Guitar World \textemdash{} Call and response phrasing (J. Smith / J. J. Nichols)}
\par Un solo es una conversación, no un catálogo de escalas. Si tocás sin escuchar lo que acabás de tocar, no hay discurso.
\begin{enumerate}
\item Tocá una pregunta corta (2 compases) que termine en tensión (2, 4 o 7).
\item Callate 2 compases. Escuchá de verdad lo que dejaste sonando.
\item Respondé con la misma idea cambiando la última nota para resolver (1 o 3).
\item Repetí 4 rondas sin repetir la respuesta. Si te quedás sin ideas, repetí la pregunta.
\end{enumerate}
\emph{Preset: escala pent-minor, 76 BPM, fondo blues.}
\subsection{10. Notas objetivo sobre la progresión}
\textbf{Armonía} \hfill \emph{fuente: MusicScene \textemdash{} Target Notes}
\par Una frase necesita destino. Elegir la nota objetivo antes de tocar (por ejemplo la 3\textordfeminine{} del acorde que viene) le da dirección armónica al solo.
\begin{enumerate}
\item Mirá la progresión: elegí una nota por acorde (3\textordfeminine{} o 7\textordfeminine{} traen el color).
\item Tocá el acorde despacio y aterrizá en esa nota, en el tiempo 1.
\item Ahora agregá lo que quieras entre un aterrizaje y el otro. Eso es un solo.
\item Cantá la nota objetivo con el acorde sonando; después encontrala en el mástil.
\end{enumerate}
\emph{Preset: escala major, 68 BPM, fondo diatonic-major.}
\subsection{11. Silencio, respiración y frases de 2 compases}
\textbf{Fraseo} \hfill \emph{fuente: Absolutely Understand Guitar \textemdash{} Scotty West (2001)}
\par El silencio es parte de la música. El aire entre frases es lo que hace que el oyente escuche lo que tocás en lugar de oír un chorro de notas.
\begin{enumerate}
\item Metrónomo a 70. Frasear en 2 compases: 2 de frase, 2 de silencio.
\item Nunca corras las frases: si una idea necesita 3 compases, esperá la vuelta.
\item Cantá tu frase antes de tocarla. Si no la podés cantar, todavía no es una idea.
\item Grabate 2 minutos y contá cuántos compases dejaste en silencio.
\end{enumerate}
\emph{Preset: escala major, 70 BPM, fondo drone.}
\subsection{12. Pulso: emocional, motriz, intelectual}
\textbf{MDM} \hfill \emph{fuente: Mauro De María \textemdash{} Composición Integral (2015)}
\par El tempo es una dimensión emocional, no solo una velocidad: \textasciitilde{}90 se percibe contemplativo, 130-140 invita a moverse, +180 sobreexcita o intelectualiza. Practicar la MISMA frase en los tres cambios de pulso es la forma más rápida de sentir para qué sirve el tempo.
\begin{enumerate}
\item Tocá una frase fija de 2 compases a 90 BPM (emocional/contemplativo).
\item Subí a 130-140: la misma frase ahora es motriz. No agregues notas.
\item Subí a 180: notá cómo el cuerpo abandona y se vuelve mental.
\item Volvé a 90 y tocá exactamente igual. La frase es la misma; la emoción, no.
\end{enumerate}
\emph{Preset: escala pent-minor, 90 BPM, fondo drone.}
\subsection{13. Ataque: legato vs. staccato}
\textbf{MDM} \hfill \emph{fuente: Mauro De María \textemdash{} Composición Integral (2015)}
\par El mismo sonido, con otro ataque, cambia la emoción: legato canta y conmueve; staccato separa, percute y se vuelve intelectual o humorístico.
\begin{enumerate}
\item Tocá un motivo legato (notas ligadas, sin cortar) y escuchá el carácter cantable.
\item Tocá el mismo motivo staccato: ahora el ritmo pasa al frente.
\item Alterná: primera frase legato, respuesta staccato.
\item Probá portato y tenuto también: separado mínimo vs. sostenido hasta el final.
\end{enumerate}
\emph{Preset: escala dorian, 66 BPM, fondo drone.}
\subsection{14. Desarrollo motívico: 10 tipos de variación}
\textbf{MDM} \hfill \emph{fuente: Mauro De María \textemdash{} Composición Integral (2015)}
\par Un solo no se construye con 20 ideas: se construye sacándole jugo a UNA. De Maria lista 10 tipos de variación; alcanzan para sostener un solo entero.
\begin{enumerate}
\item Elegí un motivo de 3 o 4 notas. Tocálo hasta que sea tuyo.
\item Aplicale variaciones: ornamental, decorativa, rítmica, simplificadora, motívica, imitativa, armónica, contrapuntística, amplificadora, estilística.
\item Por ronda, una variación distinta, sin perder de vista el motivo.
\item Ordená 4 variaciones y ya tenés un solo con forma (aparece, se complica, se resuelve).
\end{enumerate}
\emph{Preset: escala major, 72 BPM, fondo diatonic-major.}
\newpage
\section{Melod\'ia}
\subsection{Los tres sistemas mel\'odicos}
La melod\'ia es una dimensi\'on emocional propia: puede inyectar emoci\'on con independencia de la armon\'ia y del ritmo. De Mar\'ia distingue tres componentes f\'isicos —y por lo tanto tres mundos emocionales— seg\'un c\'omo se combinan tonos y semitonos.
\subsubsection{Neutral \textemdash{} escala diatónica (7 notas)}
\textbf{Normal, equilibrada}
\par Combina tonos y semitonos. Es la escala "normal" de occidente, la que nos da el parámetro de normalidad. El semitono tiene punta y dirección; el tono es más abierto y menos direccional. Sirve para hacer desear a las otras dos y para darle un descanso emocional al oyente.
\emph{Demo: escala major.}
\subsubsection{Pentafonía \textemdash{} pentatónica (sin semitonos)}
\textbf{Dulce, femenino, sin punta}
\par Sin ningún semitono: no hay nervio ni filo. Suena dulce, positiva y con mucha presencia. Es tan dulce que usada de más empalaga; alternarla con la escala neutral la hace extrañar.
\emph{Demo: escala pent-minor.}
\subsubsection{Cromatismo \textemdash{} solo semitonos}
\textbf{Nervio, punta, masculino}
\par Compuesta íntegramente por semitonos: nerviosa, filosa, agresiva. Normalmente se usa como notas de paso breves, pero sostenida mucho tiempo crea una textura totalmente distinta. Es la oposición exacta de la pentafonía.
\emph{Demo: 13 semitonos.}
\subsection{Emoci\'on de cada grado}
La misma escala mayor, escuchada desde cada grado, cambia de emoci\'on. Esta tabla es la gu\'ia m\'as r\'apida para elegir d\'onde termina una frase.
\subsubsection{C mayor}
\begin{longtable}{p{2.5cm}p{2.7cm}p{9cm}}
\textbf{Grado} & \textbf{Emoci\'on} & \textbf{Car\'acter} \\ \hline
1 \textemdash{} C & Alegre & Reposo, calma, quietud, paz \\
2 \textemdash{} Dm & Triste \ensuremath{-} & Profundo, reflexivo, optimista \\
3 \textemdash{} Em & Triste \ensuremath{-} \ensuremath{-} & Depresivo, caído, cabizbajo, muy triste \\
4 \textemdash{} F & Alegre + & Reflexivo, optimista, femenino, no efusivo \\
5 \textemdash{} G & Alegre + + & Efusivo, masculino, inquieto, activo \\
6 \textemdash{} Am & Triste & Melancólico, nostálgico, no tan triste \\
7 \textemdash{} B\textdegree{} & Raro \ensuremath{-} \ensuremath{-} \ensuremath{-} \ensuremath{-} & Nervioso, inquieto, descolocado, desencajado \\
\end{longtable}
\subsubsection{A menor natural}
\begin{longtable}{p{2.5cm}p{2.7cm}p{9cm}}
\textbf{Grado} & \textbf{Emoci\'on} & \textbf{Car\'acter} \\ \hline
1 \textemdash{} Am & Triste & Introspectivo, íntimo, melancólico, nostálgico \\
2 \textemdash{} B\textdegree{} & Raro \ensuremath{-} \ensuremath{-} \ensuremath{-} & Nervioso, inquieto, descolocado \\
3 \textemdash{} C & Alegre & Reposo, calma, quietud, paz \\
4 \textemdash{} Dm & Triste \ensuremath{-} & Profundo, mucha fuerza, cierto optimismo \\
5 \textemdash{} Em & Triste \ensuremath{-} \ensuremath{-} & Soledad, cabizbajo, angustiado, muy triste \\
6 \textemdash{} F & Alegre ++ & Femenino, dulce, suave, sin punta \\
7 \textemdash{} G & Alegre + + & Fantasía, muy optimista, positivo \\
\end{longtable}
\par \emph{Las tablas de arriba salen del cap. 18 y se verifican contra los acordes: la 3\textordfeminine{} de C es el grado más "depresivo", la 5\textordfeminine{} (dominante) el más efusivo, la 7\textordfeminine{} el más inestable. Sobre las escalas menor armónica y bachiana De María describe el mundo sonoro de forma macro \textemdash{}armónica: "mundo masculino, agresivo, con punta"; bachiana: "mundo de color, emociones más extremas y exóticas"\textemdash{} porque ahí la emoción la define el color del acorde que la sostiene.}
\subsection{Contornos mel\'odicos}
Un contorno es la forma de la frase sin importar las notas.
\begin{itemize}
\item \textbf{Arco (sube y baja)}: grados 0-1-2-3-4-3-2-1-0.
\item \textbf{Ascendente}: grados 0-1-2-3-4-5-6.
\item \textbf{Descendente}: grados 6-5-4-3-2-1-0.
\item \textbf{Pregunta (queda en tensión)}: grados 0-2-1-3.
\item \textbf{Respuesta (resuelve)}: grados 4-3-2-0.
\item \textbf{Motivo repetido}: grados 0-0-2-3-2-0.
\end{itemize}
\subsection{Ejercicio: motivo + 10 variaciones}
Un solo no se construye con veinte ideas: se construye sac\'andole jugo a una.
\begin{itemize}
\item \textbf{Ornamental}: Mismas notas, con bordados y notas de adorno alrededor.
\item \textbf{Decorativa}: Se llena el ritmo de la frase sin cambiar su dibujo.
\item \textbf{Rítmica}: Cambia la figuración: 8vos a corcheas, a tresillos, con síncopa.
\item \textbf{Simplificadora}: Al revés: menos notas, más espacios, se escucha la estructura.
\item \textbf{Motívica}: Se toma solo una célula del motivo y se explota.
\item \textbf{Imitativa}: El mismo motivo en otra cuerda, otra octava o en canon.
\item \textbf{Armónica}: Se apoya el motivo sobre otros acordes: cambia el color.
\item \textbf{Contrapuntística}: El motivo contra una segunda línea (bajo u otra voz).
\item \textbf{Amplificadora}: Se estira el motivo: duplicación de duraciones, más peso.
\item \textbf{Estilística}: El mismo motivo en otro estilo (blues, latino, barroco).
\end{itemize}
\subsection{Recursos de fraseo}
\begin{itemize}
\item \textbf{Arco melódico}: Subir, llegar a un punto climático y bajar. Sin arco, la frase no tiene forma.
\item \textbf{Densidad}: Alternar zonas densas (muchas notas) con zonas ralas. La densidad es emoción, no relleno.
\item \textbf{Registro}: Tocar la misma idea más aguda o más grave cambia la emoción sin cambiar una nota.
\item \textbf{Ataque}: Legato canta; staccato percute. Es la palanca más barata para cambiar de carácter.
\item \textbf{Intensidad}: La misma armonía en pp es dulce y tímida; en ff es heroica o violenta.
\item \textbf{Silencio}: El silencio no es ausencia de música: es la respiración de la frase.
\end{itemize}
\newpage
\section{Fuentes}
\begin{itemize}
\item \textbf{Absolutely Understand Guitar \textemdash{} Scotty West (2001)} \\ \emph{Local: absolutely-understand-guitar.pdf} \\ Lecciones 11-24: intervalos, "Must Know Scales", digitaciones, modos, pentatónicas, blues, silencio y notas de paso.
\item \textbf{Mauro De María \textemdash{} Composición Integral (2015)} \\ \emph{Local: Composicion integral.pdf} \\ Cap. 15 (tipos de variación temática) y cap. 18 (14 dimensiones: melodía, pulso, compás, ritmo, fraseo, ataque, intensidad, densidad).
\item \textbf{Total Guitarist \textemdash{} Scale Formulas} \\ \emph{Web: totalguitarist.com/lessons/theory/scales/reference/formulas/} \\ Verificación cruzada de las fórmulas por grados (mayor, menor, pentatónicas, blues, menor armónica y melódica).
\item \textbf{Premier Guitar \textemdash{} Fretboard Workshop: A Pentatonic Approach to Modes} \\ \emph{Web: premierguitar.com/lessons/fretboard-workshop-a-pentatonic-approach-to-modes} \\ Fórmulas de los 7 modos y relación pentatónica-modo.
\item \textbf{Muted.io \textemdash{} Scale Formulas, Patterns \& Intervals} \\ \emph{Web: muted.io/scale-formulas-intervals/} \\ Segunda verificación de patrones tono/semitono.
\item \textbf{Guitar World \textemdash{} Call and response phrasing (J. Smith / J. J. Nichols)} \\ \emph{Web: guitarworld.com/lessons/learn-call-and-response-soloing-with-josh-smith} \\ Ejercicio de pregunta-respuesta y desarrollo motívico ("sacarle jugo" a una idea).
\item \textbf{Improvis.io \textemdash{} Guitar Improvisation for Beginners} \\ \emph{Web: improvis.io/guides/guitar-improvisation-basics} \\ Orden de arranque: caja 1 de la pentatónica menor, fraseo antes que velocidad.
\item \textbf{MusicScene \textemdash{} Target Notes} \\ \emph{Web: musicscene.com.au/lessons/practice/improv/target-notes/} \\ Notas objetivo como destino melódico y ancla armónica.
\end{itemize}
\end{document}
